% texte.tex — f165v, vv. 252–294 (conversion initiale des travaux antérieurs)
\vers[n=252]{ΑΓ.}{Agathon}{λαμβάνετε καὶ χρῆσθ᾽· οὐ φθονῶ.}{Prenez et usez-en ; je ne m'y oppose pas.}
\vers[suite]{ΜΝ.}{Mnésiloque}{τί οὖν λάβω;}{Que faut-il prendre ?}
\vers[n=253]{ΕΥ.}{Euripide}{ὅ τι; τὸν κροκωτὸν πρῶτον ἐνδύου λαβών.}{Quoi ? prends d'abord cette crocote et revêts-la.}
\vers[n=254]{ΜΝ.}{Mnésiloque}{νὴ τὴν Ἀφροδίτην ἡδύ γ᾽ ὄζει ποσθίου.}{Par Aphrodite, elle fleure bon le priape.}
\vers[n=255]{ΑΓ.}{Agathon}{σύζωσον ἀνύσας.}{Serre-t'en vite la taille.}
\vers[suite]{ΕΥ.}{Euripide}{αἶρε νῦν στρόφιον.}{Maintenant passe-nous un strophion.}
\vers[suite]{ΑΓ.}{Agathon}{ἰδού.}{Là.}
\vers[n=256]{ΜΝ.}{Mnésiloque}{ἴθι νυν κατάστειλόν με τὰ περὶ τὼ σκέλει.}{Maintenant arrange-moi les jambes.}
\vers[n=257]{ΕΥ.}{Euripide}{κεκρυφάλου δεῖ καὶ μίτρας.}{Une résille et un bandeau ?}
\vers[suite]{ΑΓ.}{Agathon}{ἡδὶ μὲν οὖν}{Plutôt cette}
\vers[n=258]{}{}{κεφαλὴ περίθετος, ἣν ἐγὼ νύκτωρ φορῶ.}{perruque, que je porte la nuit.}
\vers[n=259]{ΕΥ.}{Euripide}{νὴ τὸν Δί᾽ ἀλλὰ κἀπιτηδεία πάνυ.}{Eh ! par Zeus, c'est même tout juste l'affaire.}
\vers[n=260]{ΜΝ.}{Mnésiloque}{ἆρ᾽ ἁρμόσει μοι;}{M'ira-t-elle ?}
\vers[suite]{ΕΥ.}{Euripide}{νὴ Δί᾽ ἀλλ᾽ ἄριστ᾽ ἔχει.}{Mais c'est parfait, par Zeus.}
\vers[n=261]{}{}{φέρ᾽ ἔγκυκλον·}{Un encycle ?}
\vers[suite]{ΑΓ.}{Agathon}{τουτὶ λάβ᾽ ἀπὸ τῆς κλινίδος.}{Prends celui-là, sur la couchette.}
\vers[n=262]{ΕΥ.}{Euripide}{ὑποδημάτων δεῖ.}{Des chaussures ?}
\vers[suite]{ΑΓ.}{Agathon}{τἀμὰ ταυτὶ λάμβανε.}{Voici les miennes, prends-les.}
\vers[n=260]{ΜΝ.}{Mnésiloque}{ἆρ᾽ ἁρμόσει μοι;}{Est-ce qu'elles m'iront ?}
\vers[suite]{ΑΓ.}{Agathon}{χαλαρὰ γοῦν χαίροις φορῶν.}{Tu préférerais en chausser de larges...}
\vers[n=264]{ΜΝ.}{Mnésiloque}{σὺ τοῦτο γίγνωσκ᾽}{Ça tu en sais quelque chose.}
\vers[suite]{ΑΓ.}{Agathon}{ἀλλ᾽ ἔχεις γὰρ ὧν δέει,}{Or çà, tu as ce qu'il te faut.}
\vers[n=265]{}{}{εἴσω τις ὡς τάχιστά μ᾽ ἐσκυκλησάτω.}{Qu'on me roule au plus vite à l'intérieur.}
\vers[n=266]{ΕΥ.}{Euripide}{ἁνὴρ μὲν ἡμῖν οὑτοσὶ καὶ δὴ γυνὴ}{Voilà notre homme devenu femme,}
\vers[n=267]{}{}{τό γ᾽ εἶδος· ἢν λαλῇς δ᾽, ὅπως τῷ φθέγματι}{du moins en apparence. Si tu parles, tâche par ton timbre de voix}
\vers[n=268]{}{}{γυναικιεῖς εὖ καὶ πιθανῶς.}{de simuler la femme, avec aisance et naturel.}
\vers[suite]{ΜΝ.}{Mnésiloque}{πειράσομαι.}{\didasc{Voix de fausset.} J'y tâcherai.}
\vers[n=269]{ΕΥ.}{Euripide}{βάδιζε τοίνυν.}{Va maintenant.}
\vers[suite]{ΜΝ.}{Mnésiloque}{μὰ τὸν Ἀπόλλω οὔκ, ἤν γε μὴ}{Non, par Apollon, à moins que}
\vers[n=270]{}{}{ὀμόσῃς ἐμοί—}{tu ne me jures...}
\vers[suite]{ΕΥ.}{Euripide}{τί χρῆμα;}{Quoi ?}
\vers[suite]{ΜΝ.}{Mnésiloque}{συσσώσειν ἐμὲ}{D'aider à me sauver}
\vers[n=271]{}{}{πάσαις τέχναις, ἤν μοί τι περιπίπτῃ κακόν.}{par toutes tes rubriques, s'il me tombe dessus quelque malheur.}
\vers[n=272]{ΕΥ.}{Euripide}{`ὄμνυμι τοίνυν αἰθέρ᾽ οἴκησιν Διός.'}{Eh bien, « je jure, par l'Éther, la demeure de Zeus ».}
\vers[n=273]{ΜΝ.}{Mnésiloque}{τί μᾶλλον ἢ τὴν Ὑποκράτους ξυνοικίαν;}{Pourquoi pas aussi bien par le logis de famille d'Hippocratès ?}
\vers[n=274]{ΕΥ.}{Euripide}{ὄμνυμι τοίνυν πάντας ἄρδην τοὺς θεούς.}{Eh bien! je jure par tous les dieux en bloc.}
\vers[n=275]{ΜΝ.}{Mnésiloque}{μέμνησο τοίνυν ταῦθ᾽, ὅτι ἡ φρὴν ὤμοσεν,}{Au surplus, souviens-toi de ceci : c'est ton cœur qui a juré,}
\vers[n=276]{}{}{ἡ γλῶττα δ᾽ οὐκ ὀμώμοκ᾽· οὐδ᾽ ὥρκωσ᾽ ἐγώ.}{et non ta langue ; je ne la tiens pas pour liée.}
\vers[n=277]{ΕΥ.}{Euripide}{ἔκσπευδε ταχέως· ὡς τὸ τῆς ἐκκλησίας}{Hâte-toi de partir. Car le signal de l'assemblée paraît sur le Thesmophorion. Moi, je m'en vais.}
\vers[n=278]{}{}{σημεῖον ἐν τῷ Θεσμοφορίῳ φαίνεται.}{(Tous sortent. Le décor du fond change et représente le Thesmophorion.}
\vers[n=279]{}{}{ἐγὼ δ᾽ ἄπειμι.}{Entrent Mnésiloque suivi d'une servante portant une ciste.)}
\vers[suite]{ΜΝ.}{Mnésiloque}{δεῦρό νυν ὦ Θρᾷτθ᾽ ἕπου.}{Par ici, Thratta, suis-moi.}
\vers[n=280]{}{}{ὦ Θρᾷττα θέασαι, καομένων τῶν λαμπάδων}{Vois, Thratta, sous la lueur fumeuse des torches ardentes,}
\vers[n=281]{}{}{ὅσον τὸ χρῆμ᾽ ἀνέρχεθ᾽ ὑπὸ τῆς λιγνύος.}{vois la foule qui monte.}
\vers[n=282]{}{}{ἀλλ᾽, ὦ περικαλλεῖ Θεσμοφόρω, δέξασθέ με}{— Thesmophores, merveilles de beauté, accueillez-moi}
\vers[n=283]{}{}{ἀγαθῇ τύχῃ καὶ δεῦρο καὶ πάλιν οἴκαδε.}{sous d'heureux auspices, ici et à mon retour au logis.}
\vers[n=284]{}{}{ὦ Θρᾷττα τὴν κίστην κάθελε, κᾆτ᾽ ἔξελε}{Thratta, dépose la ciste, et tires-en}
\vers[n=285]{}{}{τὰ πόπαν᾽, ὅπως λαβοῦσα θύσω ταῖν θεαῖν.}{le gâteau, pour que je l'offre aux deux déesses.}
\vers[n=286]{}{}{δέσποινα πολυτίμητε Δήμητερ φίλη}{Maîtresse très honorée, chère Déméter,}
\vers[n=287]{}{}{καὶ Φερρέφαττα, πολλὰ πολλάκις μέ σοι}{et toi Perséphone, puissé-je mainte et mainte fois}
\vers[n=288]{}{}{θύειν ἔχουσαν, εἰ δὲ μή, ἀλλὰ νῦν λαθεῖν}{être à même de te sacrifier. \didasc{À part :} À tout le moins, n'être pas découvert aujourd'hui.}
\vers[n=289]{}{}{καὶ τὴν θυγατέρα Χοιρίλ' ἀνδρός μοι τυχεῖν}{\didasc{Fort :} Que ma fille Chœrile trouve un mari riche}
\vers[n=290]{}{}{πλουτοῦντος, ἄλλως τ᾽ ἠλιθίου κἀβελτέρου,}{fortuné, et au surplus benêt et stupide.}
\vers[n=291]{}{}{καὶ πρὸς θυλάκιον νοῦν ἔχειν μοι καὶ φρένας.}{Et qu'elle me soit sensée et avisée sur le porte-monnaie.}
\vers[n=292]{}{}{ποῦ ποῦ καθίζωμ᾽ ἐν καλῷ, τῶν ῥητόρων}{Où donc, où m'asseoir en bonne place pour entendre les orateurs ?}
\vers[n=293]{}{}{ἵν᾽ ἐξακούω; σὺ δ᾽ ἄπιθ᾽ ὦ Θρᾷττ᾽ ἐκποδών.}{Toi, Thratta, retire-toi à l'écart,}
